\vspace{1cm}
{
\fontsize{12}{14}\selectfont

The quality of paddy seed plays a vital role in crop yield, germination rate and overall agricultural productivity. Many small and marginal farmers suffer crop loss due to substandard or counterfeit seeds that closely resemble certified products. Due to the lack of affordable and rapid verification methods, farmers rely on packaging labels and basic manual inspection, which is time consuming, inconsistent and error prone. Poor quality seeds are generally identified only after sowing, resulting in low germination rates, reduced yield and financial losses.

To address these challenges, \textbf{AgriVerify}, an AI-powered Paddy Seed Quality Assessment and Farmer Feedback System, is designed to combat the spread of counterfeit and substandard paddy seeds. The system is a web-based platform that utilizes artificial intelligence and deep learning models to automate the detection of seed quality into healthy and suspicious categories based on characteristics such as size, shape uniformity, colour, surface texture and visible defects. Farmers can scan seed samples using a simple image-based interface to receive an estimated quality assessment presented in clear, farmer-friendly language.

In addition to seed quality assessment, the system includes a Farmer Complaint and Feedback module that enables reporting of field-level issues such as poor germination, crop failure, cultivation experiences and seed performance.

The system avoids chemical analysis, specialized hardware and complex supply chain systems, making it low-cost, scalable and suitable for rural deployment. It helps farmers select better quality seeds, reduce crop failure and improve productivity. The integration of artificial intelligence with a farmer feedback mechanism makes \textbf{AgriVerify} a practical and scalable solution for modern agricultural practices and sustainable farming development.
}

